% kerr-horizons-ergosphere.tex — Kerr horizons and ergosphere cross-section
% Inputable TikZ fragment for fig:kerr-horizons-ergosphere
% Requires: tikz with calc, arrows.meta (loaded by ehbook.cls)
%
% Meridional cross-section of the Kerr black hole structure for a/M = 0.9.
% Axes: horizontal = equatorial direction (r sin theta),
%        vertical  = spin axis (r cos theta).
% In BL coordinates, surfaces of constant r are circles (spheres in 3D).
% The ergosphere boundary r_ergo(theta) is oblate, touching r+ at poles.
%
% Parameters: M = 1, a/M = 0.9.
%   r+ = 1.4359M,  r- = 0.5641M,  r_ergo(equator) = 2M.
% Scale: 1.8 cm/M.
%   r+ radius = 2.585 cm,  r- radius = 1.015 cm,  r=2M radius = 3.600 cm.
%
\begin{tikzpicture}[
    >={Stealth[length=5pt]},
    font=\footnotesize,
  ]

  % ── Ergoregion shading (between ergosphere and r+) ──────────────────
  % Even-odd rule: fill between the ergosphere path and the r+ circle.
  \fill[EHAccretionGold, opacity=0.10, even odd rule]
    plot[smooth cycle, tension=0.7] coordinates {
      ( 0.000,  2.585)
      ( 0.457,  2.593)
      ( 0.944,  2.594)
      ( 1.464,  2.535)
      ( 1.995,  2.378)
      ( 2.504,  2.101)
      ( 2.951,  1.704)
      ( 3.301,  1.201)
      ( 3.524,  0.621)
      ( 3.600,  0.000)
      ( 3.524, -0.621)
      ( 3.301, -1.201)
      ( 2.951, -1.704)
      ( 2.504, -2.101)
      ( 1.995, -2.378)
      ( 1.464, -2.535)
      ( 0.944, -2.594)
      ( 0.457, -2.593)
      ( 0.000, -2.585)
      (-0.457, -2.593)
      (-0.944, -2.594)
      (-1.464, -2.535)
      (-1.995, -2.378)
      (-2.504, -2.101)
      (-2.951, -1.704)
      (-3.301, -1.201)
      (-3.524, -0.621)
      (-3.600,  0.000)
      (-3.524,  0.621)
      (-3.301,  1.201)
      (-2.951,  1.704)
      (-2.504,  2.101)
      (-1.995,  2.378)
      (-1.464,  2.535)
      (-0.944,  2.594)
      (-0.457,  2.593)
    }
    (0, 0) circle (2.585);

  % ── Schwarzschild reference: r = 2M ─────────────────────────────────
  \draw[EHMarginGray, dotted, thin]
    (0, 0) circle (3.600);
  \node[font=\scriptsize, text=EHMarginGray, anchor=south west]
    at (2.55, 2.55) {$r = 2M$};

  % ── Ergosphere boundary ─────────────────────────────────────────────
  \draw[EHAccretionGold, thick]
    plot[smooth cycle, tension=0.7] coordinates {
      ( 0.000,  2.585)
      ( 0.457,  2.593)
      ( 0.944,  2.594)
      ( 1.464,  2.535)
      ( 1.995,  2.378)
      ( 2.504,  2.101)
      ( 2.951,  1.704)
      ( 3.301,  1.201)
      ( 3.524,  0.621)
      ( 3.600,  0.000)
      ( 3.524, -0.621)
      ( 3.301, -1.201)
      ( 2.951, -1.704)
      ( 2.504, -2.101)
      ( 1.995, -2.378)
      ( 1.464, -2.535)
      ( 0.944, -2.594)
      ( 0.457, -2.593)
      ( 0.000, -2.585)
      (-0.457, -2.593)
      (-0.944, -2.594)
      (-1.464, -2.535)
      (-1.995, -2.378)
      (-2.504, -2.101)
      (-2.951, -1.704)
      (-3.301, -1.201)
      (-3.524, -0.621)
      (-3.600,  0.000)
      (-3.524,  0.621)
      (-3.301,  1.201)
      (-2.951,  1.704)
      (-2.504,  2.101)
      (-1.995,  2.378)
      (-1.464,  2.535)
      (-0.944,  2.594)
      (-0.457,  2.593)
    };

  % ── Outer horizon r+ ────────────────────────────────────────────────
  \draw[EHHorizonCrimson, thick]
    (0, 0) circle (2.585);

  % ── Inner horizon r- ────────────────────────────────────────────────
  \draw[EHHorizonCrimson, thick, dashed]
    (0, 0) circle (1.015);

  % ── Ring singularity ────────────────────────────────────────────────
  % At r = 0, theta = pi/2: equatorial axis at origin in BL cross-section.
  \fill[EHHorizonCrimson] (0, 0) circle (2pt);

  % ── Axes ────────────────────────────────────────────────────────────
  % Spin axis (vertical)
  \draw[->, EHMarginGray, thin]
    (0, -3.2) -- (0, 3.5)
    node[above, font=\scriptsize, text=EHMarginGray] {spin axis};
  % Equatorial direction (horizontal)
  \draw[->, EHMarginGray, thin]
    (-4.1, 0) -- (4.3, 0)
    node[right, font=\scriptsize, text=EHMarginGray] {equatorial plane};

  % ── Labels ──────────────────────────────────────────────────────────

  % r+ label with leader line
  \node[font=\scriptsize, text=EHHorizonCrimson, anchor=west]
    at (2.0, -2.3) {$r_+$};
  \draw[EHHorizonCrimson, thin, shorten >=1pt]
    (2.0, -2.3) -- (1.83, -1.83);

  % r- label with leader line
  \node[font=\scriptsize, text=EHHorizonCrimson, anchor=west]
    at (1.3, -0.85) {$r_-$};
  \draw[EHHorizonCrimson, thin, shorten >=1pt]
    (1.3, -0.85) -- (0.72, -0.72);

  % r_ergo label at equatorial bulge
  \node[font=\scriptsize, text=EHAccretionGold!85!black, anchor=west]
    at (3.75, 0.35) {$r_{\text{ergo}}$};

  % Ergoregion label
  \node[font=\scriptsize, text=EHAccretionGold!85!black, rotate=55]
    at (-2.3, 1.55) {ergoregion};

  % Ring singularity label
  \node[font=\scriptsize, text=EHHorizonCrimson, anchor=north west]
    at (0.12, -0.12) {$r = 0$};

\end{tikzpicture}
